% Part: normal-modal-logic
% Chapter: completeness
% Section: introduction

\documentclass[../../../include/open-logic-section]{subfiles}

\begin{document}

\olfileid{nml}{com}{int}

\olsection{سريزه}

که $\Sigma$ يو مودالي نظام وي، نو د سموالي قضيه ثابتوي چې که
$\Sigma \Proves !A$ وي، $!A$ د مدلونو په هر هغه ټولګي
$\mClass{C}$ کښې معتبر دے چې د~$\Sigma$ د ټولو !!{formula}s
ټولې بېلګې پکې معتبرې وي۔ په ځانګړې توګه، که $\Log{K}
\Proves !A$ وي، نو $!A$ په ټولو مدلونو کښې رښتيا وي؛ که
$\Log{KT} \Proves !A$ وي، نو $!A$ په ټولو انعکاسي مدلونو کښې
رښتيا وي؛ او که $\Log{KD} \Proves !A$ وي، نو $!A$ په ټولو
مسلسلو مدلونو کښې رښتيا وي، او داسې نور۔

بشپړتيا د سموالي سرچپه لور دے: د \Log{K} بشپړتيا، د بېلګې په
توګه، دا معنا لري چې که !!a{formula}~$!A$ معتبر وي، نو
$\Proves !A$۔ د بشپړتيا ثابتول د سموالي له ثابتولو ډېر سخت دي۔
لومړے ئې نقيض عکس ته کتل ګټور دي: \Log{K} هله او يوازې هله بشپړ
دے چې هر کله $\Proves/ !A$ وي، يو مقابل مدل، يعنې داسې
مدل~$\mModel{M}$، وي چې $\mSat/{M}{!A}$۔ په همارز ډول،
(د~$!A$ په نفي کولو) دا ثابتولے شو چې هر کله
$\Proves/ \lnot !A$ وي، د~$!A$ يو مدل شته۔ د داسې مدل په جوړولو
کښې په~$!A$ کښې له شته معلوماتو ګټه اخستلے شو۔ د ځانګړو
!!{formula}s دپاره د مدلونو په موندلو کښې هم ډېر ځله همدا کوو:
د بېلګې په توګه، که د $p \lif \Box q$ دپاره مقابل مدل غواړو،
نو پوهېږو چې په هغه کښې بايد يوه نړۍ وي چې $p$ پکې رښتيا او
$\Box q$ پکې دروغ وي۔ او که $\Box q$ په يوه نړۍ کښې دروغ وي،
نو له هغې څخه د لاسرسي وړ داسې بله نړۍ بايد وي چې $q$ پکې دروغ
وي۔ همدومره پوهېدل بس دي: کومې نړۍوې کوم !!{propositional variable}s
رښتيا کوي، او کومې نړۍوې له کومو نورو نړۍو سره د لاسرسي اړيکه لري۔

خو د بشپړتيا د ثبوت په حالت کښې موږ کوم ځانګړے !!{formula}~$!A$
نۀ لرو چې مدل ورته جوړ کړو۔ غواړو وښيو چې د هر هغه $!A$ دپاره
مدل شته چې $\Proves/[\Sigma] \lnot !A$ وي۔ دا لږ تر لږه اړينه
غوښتنه ده؛ ځکه \emph{که} $\Proves[\Sigma] \lnot !A$ وي، نو د
سموالي له مخې د~$!A$ دپاره هېڅ مدل نشته (چې $\Sigma$ پکې
رښتيا وي)۔ اوس پام وکړئ چې $\Proves/[\Sigma] \lnot !A$ هله او
يوازې هله وي چې $!A$ د~$\Sigma$ په نسبت سازګار وي۔ (ياد کړئ
چې $\Sigma \Proves/[\Sigma] \lnot !A$ او
$!A \Proves/[\Sigma] \lfalse$ همارز دي۔) نو زموږ کار د~$\Sigma$
په نسبت د هر سازګار !!{formula} دپاره مدل جوړول دي۔

هغه چل چې کاروو، دا دے: د !!{formula}s داسې يو د~$\Sigma$ په نسبت
سازګار سټ پيدا کوو چې~$!A$ پکې وي، او نور فارمولونه هم ولري چې
موږ ته ووايي د~$!A$ رښتيا کوونکې نړۍ بايد څنګه وي۔ دا ډول سټونه
\emph{بشپړ} د~$\Sigma$ په نسبت سازګار سټونه دي۔ د~$!A$ د رښتيا
کولو دپاره يوازې له يوې نړۍ څخه جوړ مدل بس نۀ دے؛ مدل به څو
نړۍوې او د لاسرسي اړيکه ولري۔ هغه بشپړ د~$\Sigma$ په نسبت سازګار
سټ چې~$!A$ پکې دے، د $\Box !B$ او~$\Diamond !C$ په بڼه نور
!!{formula}s هم لري۔ $!B$ بايد په ټولو لاسرسي وړ نړۍو کښې رښتيا
وي؛ $!C$ بايد لږ تر لږه په يوه کښې رښتيا وي۔ د دې دپاره به
\emph{ټول} ممکن بشپړ د~$\Sigma$ په نسبت سازګار سټونه د نړۍو د سټ
د بنسټ په توګه واخلو۔ ستونزمنه برخه به دا وي چې وټاکو زموږ په
مدل کښې يو بشپړ د~$\Sigma$ په نسبت سازګار سټ کله له بل څخه د
لاسرسي وړ ګڼل کېږي۔

ښيو به چې په همداسې تعريف شوي مدل کښې $!A$ په يوه نړۍ کښې---چې
هغه خپله هم يو بشپړ د~$\Sigma$ په نسبت سازګار سټ دے---هله او
يوازې هله رښتيا وي چې $!A$ د هغه سټ !!a{element} وي۔ که
$!A$ د~$\Sigma$ په نسبت سازګار وي، د لږ تر لږه يوه بشپړ
د~$\Sigma$ په نسبت سازګار سټ !!a{element} به وي (دا حقيقت به
ثابت کړو)، او په دې توګه به يوه نړۍ وي چې $!A$~پکې رښتيا وي۔
نو يو داسې واحد مدل به ولرو چې د~$\Sigma$ په نسبت هر سازګار
!!{formula}~$!A$ پکې په کومه نړۍ کښې رښتيا وي۔ همدا واحد مدل
د~$\Sigma$ \emph{کانوني} مدل دے۔

\end{document}
